\section{Introduction}
\label{sec:introduction}

A bounded domain in $\C^n$ has bounded holomorphic coordinate functions
whose differentials are independent everywhere. The negatively curved
uniformization problem asks whether curvature and completeness can
force a complex manifold to possess such coordinates. In complex
dimension one, a complete simply connected surface whose curvature is
bounded above by a negative constant is biholomorphic to the disk. In higher
dimension, the existence of bounded holomorphic coordinates is a
separate global analytic requirement.

An early formulation appears in H.~Wu's work on normal families
\cite[p.~195, question~(1)]{WuNormal1967}. He asks whether a complete
simply connected K\"ahler manifold with nonpositive real sectional
curvature and holomorphic sectional curvature bounded above by a
negative constant is biholomorphic to a bounded domain. Wu and Yau
state the two-sided real-sectional version with a bounded pseudoconvex
domain as conclusion \cite[Conjecture~4.3, p.~18]{WuYauSurvey2019}:
if a simply connected complete K\"ahler manifold has every real
sectional curvature between two negative constants, must it be
biholomorphic to a bounded domain in its complex dimension?

The unit ball is already too restrictive a candidate in higher
dimension. Mostow and Siu constructed compact negatively curved
K\"ahler surfaces not covered by the ball \cite{MostowSiu1980}.
Deraux reviews their construction and obtains three-dimensional
examples \cite{Deraux2005}.
Their lifted metrics on the universal covers have two-sided negative
sectional bounds by compactness. These non-ball results, however, do
not exclude a biholomorphism to some other bounded domain; that is
the stronger obstruction sought here.

Here and throughout, sectional curvature refers to the underlying real
Riemannian metric, and completeness means geodesic completeness. A
bounded holomorphic map is a holomorphic map with bounded image in its
Euclidean target.

\begin{theorem}\label{main:theorem}
There exist a contractible domain $M\subset\C^3$, a smooth complete
K\"ahler metric $g$ on $M$, and finite constants $0<A\le B$ such that
\[
 -B\le K_g(\sigma)\le -A<0
\]
for every real two-plane $\sigma\subset T_pM$ and every $p\in M$.
No bounded holomorphic map $F:M\to\C^3$ has a nowhere-vanishing Jacobian
determinant. In particular, $M$ is not biholomorphic to a bounded domain
in $\C^3$.
\end{theorem}

Theorem~\ref{main:theorem} gives a negative answer to the stated
uniformization question, and its curvature hypotheses also place the
example within Wu's earlier question. The example has two bounded
nonconstant holomorphic functions, supplied by its base coordinates.
Its analytic obstruction concerns a fixed triple of bounded functions
with independent differentials at every point.
The domain is nevertheless Stein: the construction supplies a smooth
strictly plurisubharmonic exhaustion (Lemma~\ref{metric:complete}), so Grauert's solution
of the Levi problem applies \cite[Theorem~2 and Section~3.4]{GrauertLevi1958}.
Thus the obstruction is to bounded holomorphic coordinates, not to
holomorphic convexity.

\subsection{Intrinsic metrics and bounded functions}

Two-sided negative real sectional curvature has strong consequences for
intrinsic complex geometry. Greene and Wu established the lower
comparison for the Bergman metric \cite{GreeneWu1979}; Wu and Yau recall that result and
prove the complementary upper bound
\cite[Theorems~4 and~6]{WuYauInvariant2020}. Under completeness and
simple connectivity, Wu and Yau then compare the Bergman,
Kobayashi--Royden, complete negative K\"ahler--Einstein, and original
metrics \cite[Corollary~7]{WuYauInvariant2020}. The Bergman metric
uses square-integrable holomorphic top-degree forms. These comparisons
do not supply a fixed tuple of bounded holomorphic functions with
independent differentials everywhere. Applied to the example in
Theorem~\ref{main:theorem}, they show that all these intrinsic metrics
can be complete and mutually comparable while no such tuple exists.

The existence of even one nonconstant bounded holomorphic function is
a related question with a different conclusion. The companion paper
\cite[Theorem~1.1]{OpenAIOneSided2026} constructs, in a sufficiently
large fixed finite complex dimension, a domain diffeomorphic to
Euclidean space with a complete K\"ahler metric of real sectional
curvature at most $-1$ and only constant bounded holomorphic functions.
Its curvature is unbounded below. Here both curvature bounds are
finite and uniform, while the two bounded base coordinates remain.
The two constructions and their proofs are independent.

\subsection{Construction and proof strategy}

We construct $M$ as a family of discs over the unit ball
$\B=\{z\in\C^2:|z|<1\}$. A smooth real function $\phi$ on $\B$ determines
\[
 M_\phi=\{(z,w)\in\B\times\C:e^{\phi(z)}|w|^2<1\}.
\]
Thus the disc above $z$ has radius $e^{-\phi(z)/2}$. Put
$\psi(z)=-\log(1-|z|^2)$. Our metric is the complex Hessian of
\[
 \lambda\psi(z)-\log\bigl(1-e^{\phi(z)}|w|^2\bigr),
\]
with the real-metric convention specified in
Section~\ref{sec:reduction}. This is a specialization of Calabi's
construction from a base potential and a radial function of a Hermitian
fiber norm \cite[Section~3, equations~(3.1)--(3.2)]{Calabi1979}.
The momentum formulation of Hwang and Singer develops this framework
for general base and bundle data \cite{HwangSinger2002}.

For the logarithmic radial potential used here, Hao, Chen, and Wang
prove holomorphic and real sectional pinching results when the bundle
curvature determines the base metric
\cite[Theorems~1 and~2]{HaoChenWang2025}. Hao, Chen, Wang, and Zhang
obtain related results for complete K\"ahler--Einstein disk-bundle
metrics under Einstein and curvature hypotheses on the base
\cite[Sections~4.2--4.3]{HaoChenWangZhang2026}.
The estimate needed here concerns all real two-planes when the base
form $\lambda\partial\bar\partial\psi$ and the fiber-weight curvature
$\partial\bar\partial\phi$ are independent. We assume upper and lower
bounds for the latter form and bounds for the component derivatives
needed for curvature, expressed in centered ball coordinates. No curvature sign for the metric
$\partial\bar\partial\phi$ is assumed. We derive and estimate the
full curvature tensor under precisely these hypotheses.

The analytic task is to choose $\phi$ so that it has these derivative
bounds but cannot dominate a holomorphic real part, even after adding
the logarithmic boundary allowance forced by Cauchy's estimate for a
Jacobian. To do this, we construct radii $r_N\uparrow1$ and positive
probabilities $\mu_N$ on the unit sphere $S\subset\C^2$ such that
\[
 \int_S\left(\phi(r_N\xi)+4\log\frac1{1-r_N}\right)
                  \,d\mu_N(\xi)\longrightarrow-\infty,
\]
while $\int_S\operatorname{Re}H(r_N\xi)\,d\mu_N(\xi)
=\operatorname{Re}H(0)$ for every holomorphic $H$ on $\B$.
The two averages exclude any holomorphic real part bounded above by
the tested function plus a constant. The probabilities preserve
holomorphic evaluation at the origin, but such
measures need not preserve the logarithmic mean inequality for
holomorphic functions: Hedenmalm gives an explicit example of this
distinction \cite[Theorem~3.7]{Hedenmalm1989}.

The construction proceeds at degrees $k_j=Q^j$ for one fixed large
integer $Q$. At each degree we sum homogeneous peaks centered on a
separated set of complex lines. The projective packing and power-sum
estimate have a direct antecedent in Ryll and Wojtaszczyk
\cite[Section~2]{RyllWojtaszczyk1983}. A radial correction around each
peak annihilates every holomorphic moment exactly. Separation of the
degrees makes the corrections compatible with positivity at every
stage. Regularized logarithms built from the peak polynomials then have a
fixed negative average, which survives all subsequent corrections.
Summing these logarithms with a sufficiently large fixed coefficient
makes their negative averages dominate the boundary allowance.
Low degrees vary little
in a centered chart, while high degrees are exponentially small;
this gives the uniform derivatives through order four needed for
curvature.

The geometric task has its own delicate point. To prove negative
curvature, we must show that its negated numerator is positive on every
real two-plane. After rescaling the base directions, the leading
expression is nonnegative but can vanish on planes with nonzero base
and fiber projections. On these limiting planes the contribution of the
potentially largest error term vanishes as well. This cancellation gives stronger error
control near the degenerate planes and proves strict negativity for
one sufficiently large base parameter $\lambda$.
With $\lambda$ fixed, compactness on bounded
fiber ranges and a uniform limiting tensor at the fiber boundary give
both finite curvature bounds.

Section~\ref{sec:reduction} states the two precise ingredients and derives
Theorem~\ref{main:theorem} from them by Cauchy estimates for a putative
Jacobian. Sections~\ref{sec:peaks}--\ref{sec:densities} construct the
peaks and positive representing densities;
Sections~\ref{sec:tests}--\ref{sec:base-potential} establish their
logarithmic tests and the controlled potential.
Section~\ref{sec:hartogs} proves completeness.
Sections~\ref{sec:curvature}--\ref{sec:pinching} compute the curvature
and prove the two-sided estimate, including the limiting planes just
described.
